% =====================================================================
%  c03.tex —— §4 主问题（A∞ 增广的唯一性问题）
%  原书 p.48–50；PDF p.7 末–p.9（§5 标题之前）
%  编号公式：(7) 问题 4.1 的图式；(8) 自然性的方形图
%  定理类环境：问题 4.1、定理 4.2
%  跨块引用一律写原书字面节号与文献号 [N]，不写 \ref / \cite。
% =====================================================================

% ===== 原书 p.48 =====
\section{\texorpdfstring{$\ainf$}{A∞}增广的唯一性问题}
\label{sec:c03-4}

一般地，三角范畴$\mathcal{T}$具有$\ainf$增广时，这个增广在相差$\ainf$等价的意义下
是否唯一？这便是本文的主问题。更准确地说，考虑以下问题。

% ===== 原书 p.49 =====
\begin{problem}\label{prob:c03-4.1}
给定$\ainf$范畴$\mathcal{C}$、$\mathcal{C}'$与三角范畴等价
$\phi:\Tr(\mathcal{C})\to\Tr(\mathcal{C}')$，是否必定存在$\ainf$等价函子
$\tilde{\phi}:\Tw(\mathcal{C})\to\Tw(\mathcal{C}')$，使图式
\begin{equation}
\begin{tikzcd}[column sep=4.6em, row sep=2.8em]
  \Tw(\mathcal{C}) \arrow[r, dashed, "{\tilde{\phi}}"] \arrow[d, "H^{0}"'] &
  \Tw(\mathcal{C}) \arrow[d, "H^{0}"] \\
  \Tr(\mathcal{C}) \arrow[r, "\phi"'] & \Tr(\mathcal{C}')
\end{tikzcd}
\label{eq:c03-7}
\end{equation}
交换？
\end{problem}

\transnote{图式~\eqref{eq:c03-7} 右上角的范畴，原文印作$\Tw(\mathcal{C})$，没有撇号。
题设的函子是$\tilde{\phi}:\Tw(\mathcal{C})\to\Tw(\mathcal{C}')$，下排右端为
$\Tr(\mathcal{C}')$，按上下文此处应为$\Tw(\mathcal{C}')$。译文按原文排版，
未直接改动，读者宜按$\Tw(\mathcal{C}')$理解。}

这里，所谓该图式交换，是指合成$\phi\circ H^{0}$与$H^{0}\circ\tilde{\phi}$
（二者既不是$\ainf$函子，也不是通常意义下的函子）在与函子的自然等价相同的意义下
等价。确切地说，规定为以下内容。

\begin{itemize}[leftmargin=2em,itemsep=3pt,topsep=4pt]
  \item 对任意对象$\alpha\in\Tw(\mathcal{C})$，$\phi\circ H^{0}(\alpha)$与
        $H^{0}\circ\tilde{\phi}(\alpha)$在$\Tr(\mathcal{C}')$中同构。
        把这个同构记作
        $\theta_{\alpha}:\phi\circ H^{0}(\alpha)\to H^{0}\circ\tilde{\phi}(\alpha)$。
  \item 对任意对象$\alpha,\beta\in\Tw(\mathcal{C})$与$m_{1}^{Tw}$闭的态射
        $\Psi\in(\Tw(\mathcal{C}))^{0}(\alpha,\beta)$，以下图式交换：
\begin{equation}
\begin{tikzcd}[column sep=7em, row sep=2.8em]
  \phi\circ H^{0}(\alpha) \arrow[r, "{\phi\circ H^{0}(\Psi)}"]
  \arrow[d, "{\theta_{\alpha}}"'] &
  \phi\circ H^{0}(\beta) \arrow[d, "{\theta_{\beta}}"] \\
  H^{0}\circ\tilde{\phi}(\alpha) \arrow[r, "{H^{0}\circ\tilde{\phi}(\Psi)}"'] &
  H^{0}\circ\tilde{\phi}(\beta)
\end{tikzcd}
\label{eq:c03-8}
\end{equation}
\end{itemize}

当这样的$\tilde{\phi}$存在时，称$\tilde{\phi}$为$\phi$的提升（lift）。

\transnote{两个条件的分工：第一个只逐对象给出同构$\theta_{\alpha}$；第二个让$\Psi$
取遍$m_{1}^{Tw}$闭的零次态射，它们代表$\Tr(\mathcal{C})$中的态射，于是这个方形
恰是$\theta$对$\Tr(\mathcal{C})$中态射的自然性。所以「图式~\eqref{eq:c03-7}
交换」说的是：$\theta$给出从$\phi\circ H^{0}$到$H^{0}\circ\tilde{\phi}$的自然等价。
只是这两个复合并非真正的函子，原文才说「与自然等价同样的意义」。}

事实上，这个问题的答案是否定的。以下对象即给出反例。

\begin{theorem}[{[12, Theorem 5.5]}]\label{thm:c03-4.2}
设$\mathcal{C}$、$\mathcal{C}'$为如下极小$\ainf$范畴：
\begin{itemize}[leftmargin=2em,itemsep=3pt,topsep=4pt]
  \item 在$H(\mathcal{C})$与$H(\mathcal{C}')$中，互不相同的对象均不同构；
  \item 存在三角范畴同构$\phi:\Tr(\mathcal{C})=\Tr(\mathcal{C}')$；
  \item 这个三角范畴同构的限制给出全子范畴的同构
        $H^{0}(\mathcal{C})=H^{0}(\mathcal{C}')$。
\end{itemize}
此时，若$\mathcal{C}$与$\mathcal{C}'$并非$\ainf$同构，则$\phi$的提升
$\tilde{\phi}:\Tw(\mathcal{C})\to\Tw(\mathcal{C}')$不存在。
\end{theorem}

简单地说，若提升$\tilde{\phi}:\Tw(\mathcal{C})\to\Tw(\mathcal{C}')$存在，则其限制
应当给出$\ainf$等价函子$\mathcal{C}\to\mathcal{C}'$。而在$\mathcal{C}$与
$\mathcal{C}'$的当前条件下，$\ainf$等价函子$\mathcal{C}\to\mathcal{C}'$只要存在，
就是$\ainf$同构。由此可知，提升$\tilde{\phi}$的存在与假设矛盾。

\transnote{「极小」即$m_{1}=0$。此时$H^{0}(\mathcal{C})$与$\mathcal{C}$有相同的对象，
态射空间不必再取同调，就是原来的分次向量空间。定理~4.2 的第一条又规定$H^{0}$中
互不相同的对象不同构，于是$\ainf$等价函子在对象上只能把每个对象送到自身，在态射
空间上则给出分次向量空间的同构。这时「$\ainf$等价」与「$\ainf$同构」的差别消失，
上一段的推理正是用这一点。}

第 6 节将把满足上述定理 4.2 条件的极小$\ainf$范畴组
$(\mathcal{C},\mathcal{C}')$作为问题 4.1 的反例给出。给定$\ainf$范畴组
$(\mathcal{C},\mathcal{C}')$后，为判断三角范畴等价
$\Tr(\mathcal{C})\simeq\Tr(\mathcal{C}')$
% ===== 原书 p.50 =====
是否成立，需要考察$\Tr(\mathcal{C})$、$\Tr(\mathcal{C}')$的乘积结构
（态射的合成结构）。作为与该乘积结构对应的对象，下一章（第 5 节）首先定义
$\ainf$范畴的扭积。
